\documentclass[11pt]{article}
\usepackage{amsmath,amssymb,amsthm,mathtools}
\usepackage[margin=1in]{geometry}
\usepackage{hyperref}

\newtheorem{theorem}{Theorem}
\newtheorem{lemma}[theorem]{Lemma}
\newtheorem{proposition}[theorem]{Proposition}
\theoremstyle{definition}
\newtheorem{definition}[theorem]{Definition}
\newtheorem{remark}[theorem]{Remark}
\newtheorem{numcheck}[theorem]{Numerical Check}
\newtheorem{result}[theorem]{Gate Result}

\newcommand{\Z}{\mathbb{Z}}
\newcommand{\Q}{\mathbb{Q}}
\newcommand{\R}{\mathbb{R}}
\newcommand{\zetaK}{\zeta_{K}}
\newcommand{\Icos}{\mathcal{I}}
\renewcommand{\Theta}{\vartheta}

\title{A Closure Geometry on the $600$-Cell and Its Arithmetic Shadow\\
\large The icosian object, its $L$-function, and a certified arithmetic boundary}
\author{(VFD programme --- working draft)}
\date{2026}

\begin{document}
\maketitle

\begin{abstract}
We study a single geometric object---the icosian closure structure on the $600$-cell---
and ask what arithmetic it carries. The object is rigorous geometry: the $120$ unit
icosians, the maximal order $\Icos$ they generate over $K=\Q(\sqrt5)$, and its $E_8$
lattice. Its \emph{arithmetic shadow} is the $L$-function of its theta series, which
(by prior work of the programme) equals $\zetaK(s)\zetaK(s-1)$ \emph{exactly}, verified
by complete short-vector enumeration in the maximal icosian order. We then map the \emph{boundary} of this shadow: using a reproducible
verification gate, we certify that the object's $L$-function contains the Riemann
$\zeta$ (the unshifted factor $\zeta(s)$) only as one factor among four, and that the
proposed strong correspondence ``object $\leftrightarrow\zeta$'' is rejected. Finally we record, as clearly-marked
\emph{interpretation}, the VFD physical reading of the same object as a compactified
scale--phase system whose closure condition is a positivity witness. \textbf{No claim
is made about the Riemann Hypothesis}, and no cosmological claim is made or implied.
The mathematics is geometry-first; the physics is the frame, not the proof.
\end{abstract}

\section{The object}\label{sec:object}
This programme has been geometry-first throughout: the primitive is a shape, and its
arithmetic is discovered as a consequence. The shape here is the \emph{icosian closure
object} on the $600$-cell---the regular $4$-polytope whose $120$ vertices are the unit
icosians, i.e.\ the binary icosahedral group $2I$. These generate the icosian ring
$\Icos$, a maximal order in the definite quaternion algebra over $K=\Q(\sqrt5)$ ramified
at the two infinite places; its underlying $\Z$-lattice is a copy of $E_8$
\cite{ConwaySloane}.

Within the VFD programme this object is read physically as a compactified
scale--phase--witness system (\S\ref{sec:physics}); that reading motivates the study but
is \emph{not} used in any of the mathematics of \S\S\ref{sec:geometry}--\ref{sec:boundary}.
The companion documents \texttt{from-a-point-to-a-universe} (intuition) and
\texttt{vfd-rh-reformulation} (a falsifiable reformulation) carry the physical narrative
in full; this paper keeps it walled to \S\ref{sec:physics}.

\paragraph{What is claimed.} The geometry of \S\ref{sec:geometry}; the arithmetic shadow
of \S\ref{sec:shadow} (the exact icosian identity $L(\Theta_\Icos,s)=\zetaK(s)\zetaK(s-1)$,
recalled from \cite{IcosianTriad}); and the certified boundary of \S\ref{sec:boundary}
(the object's $L$-function contains $\zeta$ only as a factor; the strong
``object $\leftrightarrow\zeta$'' bridge candidate is rejected by an explicit gate).

\paragraph{What is not claimed.} Nothing about the location of the zeros of $\zeta$ (the
Riemann Hypothesis); no positivity, self-adjointness, operator, or arithmetic-surface
\emph{result}; no connection between any physical or cosmological model and RH. The
boundary of \S\ref{sec:boundary} concerns a specific \emph{construction}; it implies
nothing about RH.

\section{The geometry}\label{sec:geometry}
Let $\varphi=\tfrac{1+\sqrt5}{2}$, $\mathcal{O}_K=\Z[\varphi]$. The Dedekind zeta of
$K=\Q(\sqrt5)$ factors as
\begin{equation}\label{eq:dedekind}
  \zetaK(s)=\zeta(s)\,L(s,\chi_5),
\end{equation}
with $\chi_5$ the even real quadratic character mod $5$ \cite{Neukirch}. The ring
$\Icos$ is the $\Z[\varphi]$-order on the $120$ unit icosians; its $24$-element rational
sub-structure is the group of Hurwitz units (the $24$-cell / $2T$, over $\Q$), and the
remaining $96$ vertices form the $\varphi$-shell. The coordinate Galois involution
$\sigma:\sqrt5\mapsto-\sqrt5$ fixes the $24$-cell pointwise and permutes the
$\varphi$-shell. Throughout, ``the object'' (or ``substrate'') refers to $\Icos$ with
this sub-structure.

\begin{lemma}[class number one]\label{lem:h1}
The icosian maximal order $\Icos$ in the definite quaternion algebra over $K=\Q(\sqrt5)$
ramified exactly at the two infinite places has a single (left- or right-) ideal class.
\end{lemma}
\begin{proof}[Reference]
$\Icos$ is the maximal-order case with invariants $(n,d_F,D,N)=(2,5,1,1)$ in the
Kirschmer--Voight enumeration of class-number-one definite quaternion orders over
totally real fields \cite[Table 8.2]{KirschmerVoight} (see also Vign\'eras
\cite{Vigneras}). Left and right class numbers agree here by conjugation.
\end{proof}

\section{The arithmetic shadow}\label{sec:shadow}
The arithmetic carried by the object is the $L$-function of its theta series. It is
\emph{exactly} the rank-two product $\zetaK(s)\zetaK(s-1)$, verified in
\cite{IcosianTriad} by complete short-vector enumeration in the maximal icosian order.

\begin{proposition}[recalled from \cite{IcosianTriad}]\label{prop:icosian}
Let $\Theta_{\Icos}$ be the Hilbert theta series over $K$ of the norm form on the
maximal icosian order $\Icos$ (weight two). By the Siegel--Weil formula
\cite{Siegel,Weil1965} and Lemma~\ref{lem:h1} (single class), $\Theta_{\Icos}$ is the
corresponding Eisenstein series with no cuspidal contribution, and
\[
   L(\Theta_{\Icos},s)=\zetaK(s)\,\zetaK(s-1)\qquad\text{exactly.}
\]
The maximal-order representation numbers $r(\pi)=120(1+N_Q(\pi))$ at every prime
($120$ = number of unit icosians), and $r((2)^k)/120=\sigma_1((2)^k)=1,5,21,85$ at the
inert prime $(2)$, give the standard Eisenstein local factor everywhere; there is no
local correction.
\end{proposition}
\begin{remark}[on the earlier ``local-$2$ factor'']\label{rem:c2}
An earlier draft of \cite{IcosianTriad} proposed a correction $C_2(s)\ne1$ at the prime
$2$, from a count $r(2)=24$. That value is the count in the \emph{Lipschitz suborder}
$\Z[\varphi]^4$ ($8$ units), a bounded-range enumeration proxy that is non-maximal at
$2$; for the maximal icosian order $r(2)=600=120\cdot5$ and the local factor is
standard, so $C_2\equiv 1$. The correction has been withdrawn (it was also
dimensionally ill-formed: an inert prime of norm $4$ admits a local factor only in
$4^{-s}$). The identity above is exact.
\end{remark}

The two Galois components of the object are tested against the two factors of
\eqref{eq:dedekind}:
\begin{numcheck}\label{res:24cell}
The rational $24$-cell sub-object (Hurwitz, over $\Q$) is the $K=\Q$ analogue: its theta
series has on-line zeros matching those of $\zeta(s)\zeta(s-1)$ (tested to finite
precision; the same Siegel--Weil mechanism applies, the order being class-number one
over $\Q$). This is a recorded consistency check, not a theorem-grade derivation here.
\end{numcheck}
\begin{proposition}\label{prop:online}
On the line $\mathrm{Re}(s)=\tfrac12$, the zeros of $\zeta(s)\zeta(s-1)$ are exactly
those nontrivial zeros of $\zeta$ that lie on that line; $\zeta(s-1)$ contributes none.
\end{proposition}
\begin{proof}
A nontrivial zero of $\zeta(s-1)$ sits at $s=1+\rho$, $\mathrm{Re}(\rho)\in(0,1)$, hence
$\mathrm{Re}(s)\in(1,2)$; trivial zeros of $\zeta(s-1)$ are at $s=1-2n$. Neither meets
$\mathrm{Re}(s)=\tfrac12$. The statement concerns \emph{which} zeros lie on the line and
asserts nothing about whether all nontrivial zeros of $\zeta$ do.
\end{proof}
\begin{numcheck}\label{res:twin}
The $96$-vertex $\varphi$-shell (Galois-sign component) has tested on-line zeros matching
those of $L(s,\chi_5)$, with no observed coincidence with a Riemann zero in the tested
range. This too is a recorded consistency check, not a geometric derivation here.
\end{numcheck}

\section{The boundary}\label{sec:boundary}
The object's $L$-function (Proposition~\ref{prop:icosian}) has degree $4$ and, by
\eqref{eq:dedekind}, equals $\zeta(s)L(s,\chi_5)\zeta(s-1)L(s-1,\chi_5)$ up to the
local-$2$ factor. So the unshifted $\zeta(s)$ appears in it---\emph{as one factor among
four} (the shifted $\zeta(s-1)$ is also present but is irrelevant to the critical-line
count, by Proposition~\ref{prop:online}). A natural question is whether the object
isolates $\zeta$ itself. It does not, and we certify this mechanically.

\paragraph{The gate.} A proposed correspondence supplies a parameter-free procedure
producing the on-line zeros of a claimed $L$-function \emph{without taking the target as
input}. The gate checks \textbf{no-hardcoding} and \textbf{provenance} (declared-procedure
checks: it catches a \emph{declared} fitted/tuned map, relying on honest declaration,
not program analysis), \textbf{fingerprint} (nearest-neighbour spacing: fraction of
normalized gaps below $0.3$), \textbf{pointwise} (first eight zeros, tolerance $0.2$),
\textbf{density} (on-line zeros per unit height), \textbf{rigidity} (number variance),
and structural \textbf{degree}/\textbf{pole}. A hashed certificate is emitted either way;
the engine is a certifier, not a search. The geometry$\to$arithmetic enumeration itself
is the separate computation of \cite{IcosianTriad}, not re-performed by the gate.

\begin{result}[the $\zeta$-boundary]\label{res:o2}
The strong candidate ``object $\leftrightarrow\zeta$'' is \textbf{rejected}. Against
$\zeta$ (degree $1$) the object's degree-$4$ on-line content fails: \emph{degree}
($4$ vs $1$); \emph{fingerprint} ($0.077$ for the superposition vs $0.000$ for a single
$\zeta$-type spectrum); \emph{pointwise} (maximum displacement over the first eight
zeros $\approx 22.3$; first-zero displacement $\approx 7.5$); and \emph{density}
(Result~\ref{res:o2b}).
\end{result}
\begin{result}[counting route]\label{res:o2b}
A route sharing only the density diagnostic with Result~\ref{res:o2}: in the computed
on-line lists, the object's combined zero count exceeds $\zeta$'s---the count ratios are
$\approx 4.7, 3.7, 3.1$ in the windows up to heights $30,40,50$. (A leading-order
doubled density is the heuristic expectation for the two unshifted degree-one factors
visible on this line; we
do not certify the asymptotic, and present this as finite-window evidence, not a
theorem.) A list denser than $\zeta$'s on these windows is not $\zeta$'s.
\end{result}

\paragraph{Controls.}
\begin{result}[fitted-map control]\label{res:fit}
A candidate whose provenance \emph{declares} a least-squares fit to the zeros is rejected
on \textbf{provenance} alone, independent of fit quality---reproducing, mechanically, the
fitted-map trap of a prior internal audit (the ``$W5$'' audit; see
\texttt{CIRCLE\_TEST.md}).
\end{result}
\begin{result}[class-vs-individual control]\label{res:gue}
A single seeded, window-rescaled GUE operator (random Hermitian, no arithmetic; produced
by \texttt{grow\_atlas.py}) \emph{passes} the \textbf{fingerprint} check---right
universality class---but \emph{fails} \textbf{pointwise} (wrong individual spectrum); in
the saved atlas it also fails density and rigidity, though those depend on a window-scaling
convention and are not advanced as intrinsic discriminators. The load-bearing rejection is
pointwise: GUE-class membership is insufficient to be $\zeta$. This is one control instance,
not a theorem about generic GUE operators.
\end{result}

\begin{center}\small
\begin{tabular}{llp{4.2cm}}
\hline
candidate & verdict & load-bearing rejection \\
\hline
$24$-cell $\to \zeta(s)\zeta(s-1)$ (on-line) & \textsc{consistent} & --- \\
icosian $\to \zetaK(s)\zetaK(s-1)$ (exact) & \textsc{verified} & --- \\
$\varphi$-shell $\to L(s,\chi_5)$ (on-line) & \textsc{consistent} & --- \\
object $\to \zeta$ (strong bridge) & \textsc{rejected} & degree, fingerprint, pointwise, density \\
fitted map $\to \zeta$ & \textsc{rejected} & provenance \\
seeded GUE $\to \zeta$ & \textsc{rejected} & pointwise (class $\ne$ individual) \\
\hline
\end{tabular}
\end{center}
\noindent (``Consistent'' = passes the gate's declared-consistency checks against the
named target. The icosian row rests on \cite{IcosianTriad}; the $24$-cell and
$\varphi$-shell rows are finite consistency checks, not theorem-grade derivations.)

The net statement is a \emph{boundary} on the object's arithmetic shadow: it carries
$\zetaK(s)\zetaK(s-1)$ exactly, in which the unshifted $\zeta(s)$ is one factor, and it
does not---certifiably, for this candidate and this gate---isolate $\zeta$.

\section{The physical reading (interpretation)}\label{sec:physics}
\emph{This section is interpretation; nothing above depends on it, and it proves nothing.}
VFD reads the same object as a single compactified scale--phase system with a positivity
witness,
\[
   \mathcal{V}_1=\bigl([-\infty,+\infty]\times S^1,\ \pi,\ Q\bigr),\qquad
   \pi(t,\theta)=e^{t/2}e^{i\theta},
\]
in which the point ($t\to-\infty$), the unit circle ($t=0$), and infinity
($t\to+\infty$) are three boundary cuts of one object, and the critical line
$\mathrm{Re}(s)=\tfrac12$ is read as the balanced equator. The witness $Q$ is the
object's closure law, schematically $Q=H-R$ (archimedean/boundary capacity minus prime
residue), and the VFD reading of RH is the closure condition $Q[\Phi]\ge0$.

Two honest points fix the status of this reading. First, with $Q$ the Weil
explicit-formula functional, ``$Q\ge0\iff$ RH'' is the \emph{Weil criterion}---a known
theorem, restated in scale--phase coordinates, not new. Second, the only non-trivial step
would be to \emph{construct} a witness operator $A$ with $Q(f)=\langle f,Af\rangle$
geometrically from $\mathcal{V}_1$ and show $A\ge0$; that positivity \emph{is} RH and is
not claimed here. The reading is, however, \emph{falsifiable}: one searches for $f$ with
$Q(f)<0$; a single such $f$ would refute it. The companion \texttt{vfd-rh-reformulation}
develops this with the construction target and the test stated precisely. In the slogan
of the picture: \emph{geometry gives the object; arithmetic tests its closure.}

\section{Scope}\label{sec:scope}
We claim the geometry of \S\ref{sec:geometry}, the recalled arithmetic shadow of
\S\ref{sec:shadow} (with its local-$2$ factor), and the certified boundary of
\S\ref{sec:boundary}. We make no claim about the Riemann Hypothesis: the unshifted
$\zeta(s)$ enters only as one factor of the object's $L$-function, and the negative is a
statement about a
construction, not about the zeros of $\zeta$. The physical reading of
\S\ref{sec:physics} is interpretation. The gate, certificates, and atlas accompany this
paper in \texttt{vfd\_math\_engine/} (the seeded GUE control row is produced by
\texttt{grow\_atlas.py}). No proof certificate is defined or issued.

\begin{thebibliography}{9}
\bibitem{ConwaySloane} J.~H.~Conway, N.~J.~A.~Sloane, \emph{Sphere Packings, Lattices
  and Groups}, Springer.
\bibitem{Neukirch} J.~Neukirch, \emph{Algebraic Number Theory}, Springer.
\bibitem{IcosianTriad} VFD programme, \emph{The Icosian Triad on $V_{600}$}
  (\texttt{icosian-triad-v600} release, v1.1); verifies
  $L(\Theta_{\Icos},s)=\zetaK(s)\zetaK(s-1)$ exactly by complete short-vector
  enumeration in the maximal icosian order ($r(\pi)=120(1+N_Q(\pi))$ at every prime;
  $C_2=1$).
\bibitem{KirschmerVoight} M.~Kirschmer, J.~Voight, \emph{Algorithmic enumeration of
  ideal classes for quaternion orders}, SIAM J.\ Comput.\ \textbf{39} (2010), Table 8.2.
\bibitem{Vigneras} M.-F.~Vign\'eras, \emph{Arithm\'etique des alg\`ebres de
  quaternions}, Lecture Notes in Math.\ \textbf{800}, Springer (1980).
\bibitem{Siegel} C.~L.~Siegel, \emph{\"Uber die analytische Theorie der quadratischen
  Formen}, Ann.\ of Math.\ (1935).
\bibitem{Weil1965} A.~Weil, \emph{Sur la formule de Siegel\dots}, Acta Math.\ (1965).
\end{thebibliography}

\end{document}
